\documentclass[10pt]{article}

\usepackage[utf8]{inputenc}

\usepackage[T1]{fontenc}

\usepackage{amsmath}

\usepackage{amsfonts}

\usepackage{amssymb}

\usepackage[version=4]{mhchem}

\usepackage{stmaryrd}



\begin{document}

го квантования $\psi(x)$ и $\psi^{+}(x)$, удовлетворяющие перестановочным соотношениям Б.-Э. с.



\[

\psi(x) \psi^{+}\left(x^{\prime}\right)-\psi^{+}\left(x^{\prime}\right) \psi(x)=\delta\left(x-x^{\prime}\right),

\]



где $\delta\left(x-x^{\prime}\right)$ - дельта-функция. Тогда требования Б.-Э. с. оказываются выполненными и в статистич. сумме будут учитываться лишь симметричные состояния. Но в такой постановке задача вычисления статистич. суммы, а следовательно, и термодинамич. функций очень сложна и допускает приближенное решение лишь для слабо взаимодействующих систем (слабоидеальный бозе-газ) (см. [4]-[7]).



Лum.: [1] B ose S., «Z. Phys.», 1924, Bd 27, S. 384-92; [2] Эйн штей н А., Собрание научных трудов, т. 3, М., 1966 (там же в Приложении ст. Бозе Ш.- [1]); [3] Пау л и В., Релятивистская теория элементарных частиц, пер. с англ., М., 1947 (в Дополнении ст. Pauli W., «Phys. Rev.», 1940, v. 58, p. 716-22); [4] Mай е р Дж., Ге п пе рт - М а й е р М., Статистическая механика, пер. с англ., 2 изд., М., 1980; [5] Ис и х а р а А., Статистическая физика, пер. с англ., 2 изд., М., 1973; [6] Б о гол юбов Н. Н., Лекции по квантовой статистике, Иззранные труды, т. 2, Киев, 1970; [7] Л а н дау Л. Д., Л ифшиц Е. М., Статистическая физика, 3 изд., ч. 1, $\begin{array}{ll}\text { М., 1976. } & \text { Д. Н. Зубарев. }\end{array}$ Б030H - квантовая частица или квазичастица, подчиняющаяся Бозе - Эйнштейна статистике. Б.- частица с целым спином. Пространством состояний системы $N$ тождественных Б. является гильбертово пространство квадратично интегрируемых симметричных функций $N$ переменных. В случае большого канонич. ансамбля Б. описываются алгеброй операторов рождения - уничтожения, действующей в симметричном Фока пространстве (см. также Вторичное квантование).



Лит.: [1] Б е рез и н Ф. А., Метод вторичного квантования, М., 1986; [2] Боголюб о в Н. Н., Боголюбов Н. Н. (мл.), Введение в квантовую статистическую механику, М., 1984. Д. П. Санкович. БОЗОННАЯ ЛАВИНА - термин, принятый для описания начальной стадии процесса сверхизлучения, когда происходит лавинообразное (напр., экспоненциальное) нарастание числа квантов, испущенных атомной системой. Иногда термин «Б. л.» применяется как аналог термина сверхизлучение, как правило, в случаях, когда испускаемые системой частицы отличаются от квантов электромагнитного поля, напр. фотонное сверхизлучение. $\quad$ A. В. Андреев. БОЗОННОЕ БРОУНОВСКОЕ ДВИХЕНИЕ - сМ. Квантовое стохастическое исчисление.



БОЛЬЦМАНА КИНЕТИЧЕСКОЕ УРАВНЕНИЕ - ИНТеГро-дифференциальное уравнение, которому удовлетворяют неравновесные одночастичные распределения функции системы из большого числа частиц, например функция распределения $f(\boldsymbol{v}, \boldsymbol{r}, t)$ молекул газа по скоростям $\boldsymbol{v}$ и координатам $\boldsymbol{r}$, функции распределения электронов в металле, фононов в кристалле. Б. к.у.- основное уравнение микроскопич. теории неравновесных процессов (физич. кинетики), в частности кинетич. теории газов. Б. К. у. в узком смысле называют выведенное Л. Больцманом (L. Boltzmann) кинетич. уравнение для газов малой плотности, молекулы к-рых подчиняются классич. механике. Б.к.у. для квазичастиц в кристаллах, напр. для электронов в металле, называют также кинетическими уравнениями, или уравне ниями переноса.



Б. к. у. представляет собой уравнение баланса числа частиц (точнее, точек, изображающих состояние частиц) в элементе фазового объема $d \boldsymbol{v} d \boldsymbol{r}\left(\boldsymbol{v}=d v_{x} d v_{y} d v_{z} ; d \boldsymbol{r}=d x d y d z\right)$ и выражает тот факт, что изменение функции распределения частиц $f(\boldsymbol{v}, \boldsymbol{r}, t)$ со временем $t$ происходит вследствие движения частиц под действием внешних сил и столкновений между ними. Для газа, состоящего из частиц одного сорта, Б. к. у. имеет вид



\[

\frac{\partial f}{\partial t}+\left(v \frac{\partial f}{\partial \boldsymbol{r}}\right)+\frac{1}{m}\left(\boldsymbol{F} \frac{\partial f}{\partial \boldsymbol{v}}\right)=\left(\frac{\partial f}{\partial t}\right)_{\mathrm{cT}}

\]



где $\partial f / \partial t$ - изменение плотности числа частиц в элементе фазового объема $d v d r$ за единицу времени, $\boldsymbol{F}=\boldsymbol{F}(\boldsymbol{r}, t)-$ скта, действуюшая на частицу (может зависеть также и от скорости), $(\partial f / \partial t)_{\text {ст }}$ - изменение функции распределения встедствие столкновений (интеграл столкновений). Второй и третий члены уравнения (1) характеризуют соответствую- щие изменения функшии распределения в результате перемещения частиц в пространстве и действия внешних сил. Ее изменение, обусловленное столкновениями частиц, связано с уходом частиц из элемента фазового объема при так наз. прямых столкновениях и пополнением объема частицами, испытавшими «обратные» столкновения. Если рассчитывать столкновения по законам классич. механики и считать, что нет корреляции между динамич. состояниями сталкивающихся молекул, то



\[

\left(\frac{\partial f}{\partial t}\right)_{\mathrm{cT}}=\int\left(f^{\prime} f_{1}^{\prime}-f f_{1}\right) u \sigma(u, \vartheta) d \Omega d v_{1} .

\]



Здесь



\[

\begin{aligned}

f & =f(\boldsymbol{v}, \boldsymbol{r}, t), & f_{1}=f\left(\boldsymbol{v}_{1}, \boldsymbol{r}, t\right), \\

f^{\prime} & =f\left(\boldsymbol{v}^{\prime}, \boldsymbol{r}, t\right), & f_{1}^{\prime}=f\left(\boldsymbol{v}_{1}^{\prime}, \boldsymbol{r}, t\right) ;

\end{aligned}

\]



$v, v_{1}-$ скорости частиц до столкновения, $\boldsymbol{v}^{\prime}, \boldsymbol{v}_{1}^{\prime}-$ скорости тех же частиц после столкновения, $u=\left|\boldsymbol{v}-\boldsymbol{v}_{1}\right|$ - величина относительной скорости сталкивающихся частиц, $\sigma(u, \vartheta)-$ дифференциальное эффективное сечение рассеяния частиц в телесный угол $d \Omega$ в лабораторной системе координат, $\vartheta-$ угол между относительной скоростью и линией центров. Напр., для жестких упругих сфер, имеющих радиус $R$,



\[

\sigma=4 \pi R^{2} \cos \vartheta

\]



для частиц, взаимодействуюших по закону центральных сил,



\[

\sigma d \Omega=b d b d \varepsilon

\]



( $b$ - прицельный параметр, $\varepsilon$ - азимутальный угол линии центров).



Б. к.у. учитывает только парные столкновения между молекулами; оно справедливо при условии, что длина свободного пробега молекул значительно больше линейных размеров области, в к-рой происходит столкновение (для газа из упругих частиц это область порядка диаметра частиц). Поэтому Б. К. у. применимо для не слишком плотных газов. Иначе будет несправедливо основное предположение об отсутствии корреляции между состояниями сталкивающихся частиц (гипотеза молекулярного хаоса). Если система находится в статистич. равновесии, то интеграл столкновений (2) обрашается в нуль и решением Б. К. у. является Максвелла распределение.



При более строгом подходе для построения Б. к. у. исходят из Лиувилля уравнения для плотности распределения всех молекул газа в фазовом пространстве, из к-рого получают систему уравнений лля функций распределения одной, двух и т. д. молекул (Боголюбова уравнения). Эту цепочку уравнений решают с помощью разложения по степеням плотности частиц с использованием граничного условия ослабления корреляций, заменяющего гипотезу молекулярного хаоса.



Решение Б. к.у. при различных предположениях о силах взаимодействия между частицами - предмет кинетич. теории газов, к-рая позволяет вычислить кинетич. коэффициенты и получить макроскопич. уравнения для процессов переноса (вязкости, диффузии, теплопроводности).



Для квантовых газов значения эффективных сечений рассчитывают на основе квантовой механики с учетом неразличимости одинаковых частиц и того факта, что вероятность столкновения зависит не только от произведения функций распределения сталкивающихся частиц, но и от функций распределения частиц после столкновения. Для фермионов в результате этого вероятность столкновения будет уменьшаться, а для бозонов - увеличиваться. Оператор столкновения в квантовом случае принимает вид



\[

\begin{aligned}

& \left(\frac{\partial f}{\partial t}\right)_{\mathrm{cT}}=\frac{g}{h^{2}} \int\left\{f^{\prime} f_{1}^{\prime}(1 \mp f)\left(1 \mp f_{1}\right)-\right. \\

& \left.-f f_{1}\left(1 \mp f^{\prime}\right)\left(1 \mp f_{1}^{\prime}\right)\right\} u \sigma(u, \vartheta) d \Omega d p_{1},

\end{aligned}

\]



где знак минус соответствует Ферми - Дирака статистике, а знак плюс - Бозе - Эйнштейна статистике, $g$ - статистич. вес состояния ( $g=1$ для частиц со спином, равным нулю, и $g=2$ для частиц со спином $1 / 2), p-$ импульс частицы. Функции $f(p, r, t)$ нормированы так, что представляют





\end{document}